\documentclass[10pt,a4paper]{amsart}
\usepackage[margin=19mm]{geometry}
\usepackage{amsmath,amssymb,mathtools}
\usepackage[scr=rsfs,cal=euler]{mathalfa}
\usepackage{xfrac}
\usepackage[unicode,hidelinks,bookmarks=false]{hyperref}
\newcommand{\ie}{i.e.\ }
\newcommand{\cf}{cf.\ }
\newcommand{\resp}{respectively\ }
\newcommand{\Subsection}{\subsection}
\providecommand{\nobreakdash}{\nobreak\hbox{-}}
\setlength{\emergencystretch}{2em}
\multlinegap0pt
\usepackage{bm}
\usepackage{bbm}
\usepackage{dsfont}
\usepackage{xspace}
\usepackage{cancel}
\usepackage{mathtools}
\usepackage{amsthm}
\makeatletter
\@ifundefined{thm}{\newtheorem{thm}{Theorem}}{}
\@ifundefined{lem}{\newtheorem{lem}[thm]{Lemma}}{}
\makeatother
\newcommand{\N}{{\mathbb N}}
\newcommand{\R}{{\mathbb R}}
\newcommand{\cE}{{\mathcal E}}
\newcommand{\cF}{\mathcal F}
\newcommand{\cH}{{\mathcal H}}
\newcommand{\cJ}{\mathcal J}
\newcommand{\cL}{\mathcal L}
\newcommand{\di}{\diamondsuit}
\newcommand{\ee}{\varepsilon}
\newcommand{\myspan}{{\mathrm{span}}\,}
\title[Approximation of rearrangements by polarizations]{Approximation of rearrangements by polarizations}
\author{Gabriele Bianchi, Richard J. Gardner, Paolo Gronchi,, Markus Kiderlen}
\date{Source: arXiv:2509.02162v1 (CC BY 4.0)}
\begin{document}
\maketitle

\noindent\emph{Source and licence.} Approximation of rearrangements by polarizations. Version arXiv:2509.02162v1 (CC BY 4.0), distributed under the Creative Commons Attribution 4.0 International licence, as recorded in the arXiv licence metadata for this exact version and shown on the abstract page https://arxiv.org/abs/2509.02162v1. Full rights notice: licenses/arxiv-2509.02162-CC-BY-4.0.txt.\par\smallskip

\noindent\textit{Editorial scope.} Counted unit: the Thm whose source label is \texttt{thm:not\_weaklyfiniteBalls} in the article (source0.tex lines 1075--1079, source label \texttt{thm:not\_weaklyfiniteBalls}), delivered together with the article's own complete original proof of it (source0.tex lines 1081--1215, one \texttt{proof} environment of 135 lines). No mathematics is written, repaired or re-derived: statement and proof are the article's own text line for line. Required context retained uncounted: ball (lines 364--366); eq:psi1 (lines 380--382); eq:omregn (lines 391--395); lem\_balls\_close (lines 1004--1010); lem\_almost\_affine (lines 1030--1036). Every definition, display and label of those lines is retained unchanged. Omitted from this package and not redistributed: 1--363, 367--379, 383--390, 396--1003, 1011--1029, 1037--1074, 1080--1080, 1216--1554. Nothing inside a retained excerpt is omitted or altered: every retained body is delivered exactly as the article's lines are written, so no omission receipt of any kind accompanies this package. Citations: the retained excerpts carry no citation command at all, so no bibliography entry of the article is transcribed and no reference is redistributed. Reference adaptation: none was needed and none was made; every reference inside the delivered unit resolves inside it, so no unresolved reference or citation is produced. Printed slips kept literal: no printed word, symbol, hypothesis or number of the counted statement or of its proof was altered. Layout and class: the article's own class is \texttt{amsart} with packages including amssymb, etoolbox, graphicx, stmaryrd, color, pifont, enumerate; the wrapper is the lane's pinned \texttt{amsart} editorial wrapper at 10pt on A4, which restates the theorem-like environments the retained text needs and adds the article's own macro definitions that the retained text needs (\texttt{\textbackslash N}, \texttt{\textbackslash R}, \texttt{\textbackslash cE}, \texttt{\textbackslash cF}, \texttt{\textbackslash cH}, \texttt{\textbackslash cJ}, \texttt{\textbackslash cL}, \texttt{\textbackslash di}, \texttt{\textbackslash ee}, \texttt{\textbackslash myspan}), with extra wrapper packages (bm, bbm, dsfont, xspace, cancel, mathtools). The delivered numbering is the unit's own. Assets: no retained span contains a figure environment, a graphics-inclusion command, a PDF-page inclusion, a table or a TikZ picture; no image file is copied into this package, the package declares no asset dependency and no image file is redistributed with it. Licence: version arXiv:2509.02162v1 of this article is distributed under the Creative Commons Attribution 4.0 International licence, as recorded in the arXiv licence metadata for this exact version and shown on the abstract page https://arxiv.org/abs/2509.02162v1. A full rights notice is delivered at licenses/arxiv-2509.02162-CC-BY-4.0.txt. Characters: every retained excerpt is pure ASCII, so the delivered engine, XeLaTeX with its default Latin Modern text font, reports no missing character.
\begin{equation}\label{ball}
\di B(x+tu,r)=B\left(x+\varphi_{\di}(t)u,r\right),
\end{equation}

\begin{align}\label{eq:psi1}
\di B(x,r)=B(\psi_{\di,r}(x),r)
\end{align}

\begin{equation}\label{eq:omregn}
\psi_{\di,r}(x)=x|u^\perp+\varphi_{\di}(\langle x,u\rangle)u
\quad{\text{and}}~~\quad
\varphi_\di(t)=\langle \psi_\di(tu),u\rangle
\end{equation}

\begin{lem}\label{lem_balls_close}
Let $\cE\subset \cL^n$ be a class that contains all balls, and let  $\di,\heartsuit$ be maps in $\cJ(\cE)$ with associated contractions $\psi_{\di,r},\psi_{\heartsuit,r}$, respectively, from \eqref{eq:psi1}.
For all $x\in \R^n$ and $r>0$ with $(\di B(x,r))\cap \heartsuit B(x,r)\ne\emptyset$, we have
\[
\|\psi_{\di,r}(x)-\psi_{\heartsuit,r}(x)\|\le \frac{n}{2r^{n-1}\kappa_{n-1}}\|1_{\di B(x,r)}-1_{\heartsuit B(x,r)}\|_1.
\]
\end{lem}

\begin{lem}\label{lem_almost_affine}
Let $\psi:\R^n\to \R^n$ be a contraction and let $\ee\ge 0$. If $x,x'\in \R^n$ are such that $L=\|\psi(x)-\psi(x')\|$ satisfies $\|x-x'\|\le L+\ee$, then
\[
\|\psi((1-t)x+tx')-((1-t)\psi(x)+t\psi(x'))\|
\le \sqrt{\ee(2L+\ee)}, \qquad 0\le t\le 1.
\]
\end{lem}

\begin{thm}\label{thm:not_weaklyfiniteBalls}
Let $H$ be a hyperplane in $\R^n$ and let $\di\in \cJ_H(\cF_n^n)$. Furthermore, let $\cJ'\subset \cJ(\cF_n^n)$ be such that for each map in $\cJ'$, the associated contraction from \eqref{eq:psi1} is independent of $r>0$.

If $\di$ is weakly approximable on ${\cF}_n^n$ by maps in $\cJ'$, then there is a sequence $(\di_k)$ from $\cJ'$ such that the associated contractions $\psi_{\di_k}$ converge uniformly on $B^n$ to the contraction $\psi_{\di}$ associated with $\di$.
\end{thm}

\begin{proof}
We shall ignore sets of measure zero in the proof. Let $\ee>0$. It will suffice to show that there is a map $\heartsuit\in\cJ'$ such that its associated contraction $\psi_{\heartsuit}$ from \eqref{eq:psi1} satisfies $\|\psi_{\heartsuit}(x)-\psi_\di(x)\|\le \ee$ for all $x\in B^n$.

Since $\di$ respects $H$-cylinders, we have $(\di B^n)|H=
B(\psi_\di(o),1)|H\subset B^n+H^{\perp}$, so $\psi_\di(o)\in H^\perp$. By applying a translation by $-\psi_\di(o)$ to the entire construction, if necessary, we may assume that $\psi_\di(o)=o$. We may also assume that $H=e_n^\perp$.  Then from the second relation in \eqref{eq:omregn}, we have $\varphi_{\di}(0)=\langle \psi_{\di}(o),e_n\rangle=\langle o,e_n\rangle=0$, where $\varphi_{\di}$ is the contraction associated with $\di$ from \eqref{ball}.

We start by constructing a set $K\in \cF_n^n$ to which the weak approximation property will be applied. Let $2\le m\in \N$ satisfy $\sqrt{n-1}/(2m)\le \ee/4$, and consider the family $\{R_1,\ldots,R_{m'}\}$ of $m'=2(2m+1)^{n-1}$ rays
\[
\{(t,i_2/m,\ldots,i_n/m):t\ge 0\},\qquad
\{(t,i_2/m,\ldots,i_n/m):t\le 0\},
\]
$i_2,\ldots,i_n=-m,\dots,m$, each emanating from the hyperplane $e_1^{\perp}$ and parallel to $e_1$.  The order of enumeration is irrelevant here, apart from the convention that $R_{2i-1}\cup R_{2i}$ is a line for $i=1,\dots,m'/2$. For each ray, we now define iteratively a ball with center on that ray.

Let $B_{1}=B(x_1,r_1)$, where $x_1\in R_1$, $r_1=1+1/m'$, and where $B_{1}$ and the cube $[-1,1]^n$ meet only on their boundaries.  Suppose that $i=1,\dots,m'-1$ and we have defined balls $B_j$, $j=1,\dots,i$ such that $B_j$ has its center on $R_j$.  Let $C_i=\rho_i B^n+\myspan\{e_n\}$ be the smallest cylinder with axis parallel to $e_n$ containing the balls $B_1,\dots,B_i$, and let $B_{i+1}=B(x_{i+1},r_{i+1})$, where $x_{i+1}\in R_{i+1}$, $r_{i+1}=1+(i+1)/m'$, and where $B_{i+1}$ and the enlarged cylinder $C_i+4B^n$ meet only on their boundaries.

Let $C=\rho B^n+\myspan\{e_n\}$ be the smallest cylinder with axis parallel to $e_n$ containing the balls $B_1,\dots,B_{m'}$. Let $S_1,\dots,S_{2(n-1)}$ be the rays emanating from the origin and spanned by $\pm e_1,\dots,\pm e_{n-1}$, and for $j=1,\dots,2(n-1)$, let $D_j=B(y_j,(n-1)^{-1/2})$, where $y_j\in S_j$ and where $D_j$ and the enlarged cylinder $C+4B^n$ meet only on their boundaries.

Define
\[
K_0= {\textstyle \bigcup}_{j=1}^{2(n-1)} D_j ~~\quad{\text{and}}~~\quad
K=K_0\cup {\textstyle \bigcup}_{i=1}^{m'} B_i.
\]
The set $K$ is a finite disjoint union of balls with radii in $\{(n-1)^{-1/2}\}\cup[1,2]$. It is easily checked that by construction, if $K_1$ and $K_2$ are different components (balls) of $K$, then
\begin{equation}\label{eq:disjointn}
d(K_1|H,K_2|H)\ge 4,
\end{equation}
where $d(A,B)=\inf\{\|a-b\|: a\in A, b\in B\}$ is the usual distance between sets. Let $B_c=B(o,\rho_c)$ be the circumball of $K$ and note that all the balls in $K$ are contained in its interior, with the exception of the balls $D_j$, and that $D_j\cap\partial B(o,\rho_c)=S_j\cap \partial B(o,\rho_c)$, $j=1,\dots,2(n-1)$.

Next, we determine $\di K$. Since $\di\in \cJ_H(\cF_n^n)$ and $H=e_n^{\perp}$, relation
\eqref{eq:omregn} shows that if $x\in e_n^{\perp}$, then $\psi_{\di}(x)=x+\varphi_{\di}(0)e_n=x$. It follows from \eqref{eq:psi1} that $\di D_j=D_j$, $j=1,\dots, 2(n-1)$. Therefore $D_j=\di D_j\subset \di K_0$ by monotonicity, which gives $K_0\subset \di K_0$ and hence $\di K_0=K_0$ by the measure-preserving property of $\di$.

If $B=B(x,r)$ is a component of $K$, then $\di B\subset (B|H)+H^\perp$ as $\di$ respects $H$-cylinders, so \eqref{eq:disjointn} yields
\begin{equation}\label{eq:disjointn_din}
d\big((\di K_1)|H,(\di K_2)|H\big)\ge 4
\end{equation}
for different components $K_1, K_2$ of $K$. Consequently, the images under $\di$ of the components of $K$ are disjoint balls. In particular, \eqref{eq:psi1} and \eqref{eq:omregn} imply that the components of $\di \left(\cup_{i=1}^{m'} B_i\right)$ are the disjoint balls
$B\big(x_i^\di,r_i\big)$, with
\[
x_i^\di=x_i+\left(\varphi_\di(\langle x_i,e_n\rangle)-\langle x_i,e_n\rangle\right)e_n
\]
for $i=1,\dots,m'$, and $\di K=K_0\cup\cup_{i=1}^{m'} B(x_i^\di,r_i).$

Choose $\ee_1>0$ such that
\begin{equation}\label{eqepchoose}
\ee_1<\min\left\{1,\frac{{\cH}^n(D_1)}{2}, \frac{5\kappa_{n-1}}{n}, \frac{n\kappa_n}{m'}, \frac{(\kappa_{n-1}\ee)^2}{50n^2\rho_c}\right\}.
\end{equation}
(The list of quantities in the minimum is for convenience; in fact, one can show that the second quantity is less than the third.) By assumption, there is a map $\heartsuit\in \cJ'$ such that
$\|1_{\di K}-1_{\heartsuit K}\|_1\le \ee_1$.

The circumball $B_c=B(o,\rho_c)$ of $K$ has its center at the origin, so $\psi_\di(o)=o$ gives $\di B_c=B_c$. Since $\heartsuit K\subset \heartsuit B_c=B(\psi_{\heartsuit}(o),\rho_c)$ we obtain
\begin{align*}
\cH^n\left(D_j\setminus B(\psi_{\heartsuit}(o),\rho_c)\right)
&=\cH^n\left((\di D_j)\setminus \heartsuit B_c\right)\\
&\le \cH^n((\di K)\setminus \heartsuit K)\le \|1_{\di K}-1_{\heartsuit K}\|_1\le \ee_1<{\cH}^n(D_j)/2
\end{align*}
for $j=1,\dots,2(n-1)$.  It follows that $B(\psi_{\heartsuit}(o),\rho_c)$ contains the centers of all the balls in $K_0$, that is, the points $\pm (\rho_c-(n-1)^{-1/2}) e_j$, for $j=1,\dots,n-1$.  Thus each vector $\psi_{\heartsuit}(o)\pm (\rho_c-(n-1)^{-1/2}) e_j$, $j=1,\dots,n-1$, has Euclidean norm less than or equal to $\rho_c$.  By considering their projections on the $j$th coordinate axis, we see that
$$|\langle \psi_{\heartsuit}(o), e_j\rangle|\le \frac1{\sqrt{n-1}}$$
for $j=1,\dots,n-1$.  From this and $\psi_{\heartsuit}(o)|H=\sum_{j=1}^{n-1}\langle \psi_{\heartsuit}(o),e_j\rangle e_j$, we obtain
\begin{equation}\label{eq:circumcentern}
\|\psi_{\heartsuit}(o)|H\|\le 1.
\end{equation}
Choose $i\in \{1,\ldots, m'\}$. The ball $\heartsuit B_i$ must meet the interior of at least one component of $\di K$, as otherwise, we would have the contradiction
\begin{eqnarray*}
{\cH}^n(B_i) = {\cH}^n(\heartsuit B_i)&=&\cH^n\left( (\heartsuit B_i)\setminus \di K\right)\\
&\le& \cH^n\left( (\heartsuit K)\setminus \di K\right)
\le \|1_{\di K}-1_{\heartsuit K}\|_1\le
\ee_1< {\cH}^n(D_1)/2\le {\cH}^n(B_i)/2.
\end{eqnarray*}

We claim that $\di B_i$ is the unique component of $\di K$ that meets $\heartsuit B_i$. To see this, first recall that $B_i$ is contained in the cylinder $C_i=\rho_i B^n+\myspan\{e_n\}$ and that $|\langle x_i,e_n\rangle|\le 1$. This gives $B_i\subset B(o,\rho_i+1)$, since
\[
\|x_i\|+r_i\le \sqrt{(\rho_i-r_i)^2+1}+r_i\le (\rho_i-r_i)+1+r_i=\rho_i+1.
\]
Using \eqref{eq:circumcentern}, we obtain
\begin{align*}
\heartsuit B_i&\subset
\heartsuit B\left(o,\rho_i+1\right)
=B\left(\psi_{\heartsuit}(o),\rho_i+1\right)\subset
(\rho_i+2) B^n+\myspan\{e_n\}=C_i+2B^n.
\end{align*}
By the construction of $K$, the only components of $K$ contained in the cylinder $C_i+2B^n$ are $B_1,\ldots,B_i$, and since $\di$ respects $H$-cylinders, the only components of $\di K$ contained in  $C_i+2B^n$ are  $\di B_1,\ldots,\di B_i$.

Because $\heartsuit B_i$ has radius at most 2, by \eqref{eq:disjointn_din} it meets the interior of exactly one component of $(C_i+2B^n)\cap \di K$.  Let $\di B'$ be this component, and suppose that $B'\neq B_i$. Then our construction ensures that the radius $r$ of $B'$ satisfies $r_i-r\ge 1/m'$.  Thus, using $r_i,r\ge 1$ and $(\heartsuit B_i)\cap \di K=(\heartsuit B_i)\cap \di B'$, we find that
\begin{align*}
\frac{n\kappa_n}{m'}&\le  \kappa_n(r_i-r)\sum_{k=0}^{n-1}r_i^kr^{n-1-k}= \kappa_n(r_i^n-r^n)
=\cH^n(\heartsuit B_i)-\cH^n(\di B')
\le \cH^n\left( (\heartsuit B_i)\setminus\di B'\right)\nonumber\\
& = \cH^n((\heartsuit B_i)\setminus \di K)\le \cH^n\left( (\heartsuit K)\setminus\di K\right)\le \|1_{\di K}-1_{\heartsuit K}\|_1\le \ee_1,
\end{align*}
contradicting the choice \eqref{eqepchoose} of $\ee_1$. Therefore, $(\di B_i)\cap \heartsuit B_i\neq\emptyset$, as claimed.

Since $\di B_i$ meets $\heartsuit B_i$, Lemma~\ref{lem_balls_close} with $x=x_i$ and $r=r_i$ implies that
\begin{equation}\label{first_fixed_pointn}
\left\|\psi_{\di}(x_i)-\psi_{\heartsuit}(x_i)\right\|
\le \frac{n}{2r_i^{n-1}\kappa_{n-1}}\|1_{\di B_i}-1_{\heartsuit B_i}\|_1
\le\frac{n}{2\kappa_{n-1}} \|1_{\di K}-1_{\heartsuit K}\|_1\le \frac{n\ee_1}{2\kappa_{n-1}}.
\end{equation}

For $i=1,\dots,m'/2$, the points $x=x_{2i-1}$ and $x'=x_{2i}$ belong to the line $\ell_i=R_{2i-1}\cup R_{2i}$, which is parallel to $H$. Since $\di$ respects $H$-cylinders, \eqref{eq:omregn} implies that $\|\psi_\di(x)-\psi_\di(x')\|=\|x-x'\|$, so $\psi_\di$ is affine when restricted to $[x,x']$.  The last equality and \eqref{first_fixed_pointn} yield
\begin{equation}\label{eq:rho-1}
\|\psi_{\heartsuit}(x)-\psi_{\heartsuit}(x')\|\ge \|x-x'\|-\frac{n\ee_1}{\kappa_{n-1}}.
\end{equation}
Let $z=(1-t)x+tx'$, where $0<t<1$.  Then \eqref{eq:rho-1} allows us to apply Lemma~\ref{lem_almost_affine} with $\ee=n\ee_1/\kappa_{n-1}$ and $\psi=\psi_{\heartsuit}$, and \eqref{first_fixed_pointn}, to obtain
\begin{eqnarray}
\left\|	\psi_{\heartsuit}(z)-\psi_{\di}(z)\right\|
&=& \left\|	\psi_{\heartsuit}(z)-((1-t)\psi_{\di}(x)+t\psi_{\di}(x'))\right\|\nonumber\\
&\le& \left\|\psi_{\heartsuit}(z)-((1-t)\psi_{\heartsuit}(x)
+t\psi_{\heartsuit}(x'))\right\|+\nonumber\\
& &+(1-t)\left\|\psi_{\heartsuit}(x)-\psi_{\di}(x)\right\|+
t\left\|\psi_{\heartsuit}(x')-\psi_{\di}(x')\right\|\nonumber\\
&\le &\sqrt{\frac{n\ee_1}{\kappa_{n-1}}\left(2L+\frac{n\ee_1}{\kappa_{n-1}}
\right)}+\frac{n\ee_1}{2\kappa_{n-1}}.
\label{atlastn}
\end{eqnarray}
By \eqref{eq:disjointn}, we have $\|x-x'\|\ge 6$, and the choice \eqref{eqepchoose} of $\ee_1$ gives $n\ee_1/\kappa_{n-1}\le 5$.  These together imply the second inequality in the following chain:
\[
\ee_1\le 1\le \|x-x'\|-\frac{n\ee_1}{\kappa_{n-1}}\le L=\|\psi_{\heartsuit}(x)-\psi_{\heartsuit}(x')\|\le \|x-x'\|\le 2\rho_c.
\]
The other relations follow from the choice \eqref{eqepchoose} of $\ee_1$, \eqref{eq:rho-1}, and the fact that $\psi_\heartsuit$ is a contraction.  It is straightforward to check that for integer $n\ge 2$, we have $n/\kappa_{n-1}\ge 3/\pi$, the value when $n=3$ or 4, and hence $2\le (3n)/\kappa_{n-1}$.  All these estimates,  \eqref{atlastn}, and the choice \eqref{eqepchoose} of $\ee_1$ yield
\begin{equation}\label{eq6425}\left\|	\psi_{\heartsuit}(z)-\psi_{\di}(z)\right\|\le\sqrt{\frac{n\ee_1}{\kappa_{n-1}}\left(\frac{3nL}{\kappa_{n-1}}
+\frac{nL}{\kappa_{n-1}}\right)}+\frac{n\sqrt{L\ee_1}}{2\kappa_{n-1}}
= \frac{5n\sqrt{L\ee_1}}{2\kappa_{n-1}}\le \frac{5n\sqrt{2\rho_c\ee_1}}{2\kappa_{n-1}}\le\frac{\ee}{2}.
\end{equation}

By construction, $[-1,1]^n\cap \ell_i\subset [x, x']$, so \eqref{eq6425} holds for all $z\in [-1,1]^n\cap \cup_{i=1}^{m'/2}\ell_i$. Now for any $w\in B^n$, there is an $i\in\{1,\dots,m'/2\}$ and a $z\in [-1,1]^n\cap\ell_i$ with $\|w-z\|\le \sqrt{n-1}/(2m)\le \ee/4$.  Since $\psi_{\heartsuit}$ and $\psi_\di$ are contractions, \eqref{eq6425} gives
\begin{align*}
\left\|	\psi_{\heartsuit}(w)-\psi_\di (w)\right\|
&\le \left\|\psi_{\heartsuit}(w)-\psi_{\heartsuit} (z)\right\|
+\left\|\psi_{\heartsuit} (z)-\psi_{\di}(z)\right\|+
\left\|\psi_{\di}(z)-\psi_\di (w)\right\|\\
&\le \left\|\psi_{\heartsuit}(z)-\psi_\di (z)\right\|
+2\left\|w-z\right\|\le \ee,
\end{align*}
as required.
\end{proof}

\end{document}
